\documentclass[10pt]{article}
\usepackage[utf8]{inputenc}
\usepackage[T1]{fontenc}
\usepackage{amsmath}
\usepackage{amsfonts}
\usepackage{amssymb}
\usepackage[version=4]{mhchem}
\usepackage{stmaryrd}

\begin{document}
\begin{enumerate}
  \item Solve the equation $2 \tan ^{2} \theta-6 \sec \theta+10=4$. Provide simplified general solution(s). Then, provide all solutions on the interval ( $-\pi, 0]$. (6 points)
\end{enumerate}

General Solution(s): $\_\_\_\_$\\
On $(-\pi, 0]$ : $\_\_\_\_$\\
2. Find all solutions to the equation $\sin ^{2}(3 x)-\sin ^{2}(x)=\cos ^{2}(x)$. Provide unique general solutions. Then, provide all solutions on the interval $[0,2 \pi)$. ( 5 points)

General Solution(s): $\_\_\_\_$\\
On $[0,2 \pi)$ : $\_\_\_\_$\\
3. Verify the identity $\frac{\sec (\theta)-\cos (\theta)}{\sin (\theta)}=\tan (\theta)$ (3 points)\\
4. Write $\frac{2-2 \cos ^{2}(x)+\sin (2 x)}{2 \sin (x)}$ as the sum or difference of two separate trig ratios. (4 points)\\
5. Write $\frac{\cos (x)+\sin ^{3}(x) \cos (x)}{\cos (x) \sin (x)}+\cos ^{2}(x)-1$ as a single trig ratio. (4 points)

\section*{Unit 9 Assessment}
\begin{enumerate}
  \setcounter{enumi}{5}
  \item Solve $2 \sin ^{2}(x) \cos ^{2}(x)-\cos ^{2}(x)=2 \sin ^{2}(x)-1$. Provide simplified general solution(s). Then, provide all solutions on the interval $\left(-\pi, \frac{\pi}{2}\right)$. (4 points)
\end{enumerate}

General Solution(s): $\_\_\_\_$\\
On $\left(-\pi, \frac{\pi}{2}\right)$ : $\_\_\_\_$\\
6. If $\sin (\beta)=\frac{2}{7}$ and $\cos (\alpha)=-\frac{3}{5}$ for $\frac{\pi}{2}<\beta<\pi<\alpha<2 \pi$, find each of the following (1 point each)

\begin{center}
\begin{tabular}{|l|l|}
\hline
$\sin (\alpha-\beta)$ & $\sin (\beta+\alpha)$ \\
\hline
$\cos (\beta-\alpha)$ & $\cot \left(\frac{\beta}{2}\right)$ \\
\hline
$\tan (\beta-\alpha)$ & $\sec (2 \beta)$ \\
\hline
\end{tabular}
\end{center}

\begin{enumerate}
  \item Solve the equation $3 \cot ^{2} \theta+4 \csc \theta+4=5$. Provide simplified general solution(s). Then, provide all solutions on the interval ( $-\pi, 0]$. (6 points)
\end{enumerate}

General Solution(s): $\_\_\_\_$\\
On $(-\pi, 0]$ : $\_\_\_\_$\\
2. Find all solutions to the equation $\cot ^{2}(3 x)-\cos ^{2}(x)=\sin ^{2}(x)$. Provide unique general solutions. Then, provide all solutions on the interval $[0, \pi)$. ( 5 points)

General Solution(s): $\_\_\_\_$\\
On $[0, \pi)$ : $\_\_\_\_$\\
3. Verify the identity $\frac{\tan (\theta)}{\sec (\theta)}+\frac{\cos ^{2}(\theta)}{1+\sin (\theta)}=1$ (3 points)\\
4. Write $\frac{2-2 \cos ^{2}(x)+\sin (2 x)}{2 \sin (x)}$ as the sum or difference of two separate trig ratios. (4 points)\\
5. Write $\frac{\sin (x)+\sin (x) \cos (x)}{\sin (x) \cos (x)}-\frac{\cos (x) \csc (x)+\sin (x) \sec (x)}{\sec (x) \csc (x)}$ as a single trig ratio. (4 points)

\section*{Unit 9 Assessment}
\begin{enumerate}
  \setcounter{enumi}{5}
  \item Solve $4 \sin ^{2}(x) \cos ^{2}(x)-4 \cos ^{2}(x)=3 \sin ^{2}(x)-3$. Provide simplified general solution(s). Then, provide all solutions on the interval $\left(-\pi, \frac{\pi}{2}\right)$. (4 points)
\end{enumerate}

General Solution(s): $\_\_\_\_$\\
On $\left(-\pi, \frac{\pi}{2}\right)$ : $\_\_\_\_$\\
6. If $\sin (\beta)=\frac{4}{7}$ and $\cos (\alpha)=-\frac{4}{5}$ for $\frac{\pi}{2}<\beta<\pi<\alpha<2 \pi$, find each of the following (1 point each)

\begin{center}
\begin{tabular}{|l|l|}
\hline
$\sin (\alpha+\beta)$ & $\sin (\beta-\alpha)$ \\
\hline
$\cos (\beta-\alpha)$ & sec $\left(\frac{\beta}{2}\right)$ \\
\hline
$\tan (\beta+\alpha)$ & $\csc (2 \alpha)$ \\
\hline
\end{tabular}
\end{center}

\begin{enumerate}
  \item Solve the equation $5 \tan ^{2} \theta-\sec \theta+10=9$. Provide simplified general solution(s). Then, provide all solutions on the interval ( $-\pi, 0$ ). (6 points)
\end{enumerate}

General Solution(s): $\_\_\_\_$\\
On ( $-\pi, 0$ ): $\_\_\_\_$\\
2. Find all solutions to the equation $\cos ^{2}(3 x)-\sin ^{2}(x)=\cos ^{2}(x)$. Provide unique general solutions. Then, provide all solutions on the interval $[0,2 \pi)$. ( 5 points)

General Solution(s): $\_\_\_\_$\\
On $[0,2 \pi)$ : $\_\_\_\_$\\
3. Verify the identity $\frac{\sin (\theta)}{\sec (\theta)-\cos (\theta)}=\cot (\theta)$ (3 points)\\
4. Write $\frac{2 \cos ^{3}(x)}{\sin (2 x)}$ as the sum or difference of two separate trig ratios. (4 points)\\
5. Write $\frac{2 \sin (x) \csc (2 x)+\sin ^{2}(2 x)}{\sin (2 x) \csc (2 x)}+\frac{\cos ^{4}(2 x)}{1-\sin ^{2}(2 x)}-1$ as a single trig ratio. (4 points)

\section*{Unit 9 Assessment}
\begin{enumerate}
  \setcounter{enumi}{5}
  \item Solve $2 \sin ^{2}(x) \cos ^{2}(x)+2 \sin ^{2}(x)=\cos ^{2}(x)+1$. Provide simplified general solution(s). Then, provide all solutions on the interval $\left(-\pi, \frac{\pi}{2}\right)$. (4 points)
\end{enumerate}

General Solution(s): $\_\_\_\_$\\
On $\left(-\pi, \frac{\pi}{2}\right)$ : $\_\_\_\_$\\
6. If $\sin (\beta)=\frac{5}{13}$ and $\cos (\alpha)=-\frac{2}{5}$ for $\frac{\pi}{2}<\beta<\pi<\alpha<2 \pi$, find each of the following (1 point each)

\begin{center}
\begin{tabular}{|l|l|}
\hline
$\sin (\alpha-\beta)$ & $\sin (\beta+\alpha)$ \\
\hline
$\cos (\beta+\alpha)$ & $\tan \left(\frac{\beta}{2}\right)$ \\
\hline
$\tan (\beta-\alpha)$ & $\csc (2 \alpha)$ \\
\hline
\end{tabular}
\end{center}


\end{document}